% Bagian: computability
% Bab: recursive-functions
% Subbagian: partial-functions

\documentclass[../../../include/open-logic-section]{subfiles}

\begin{document}

\olfileid[id]{cmp}{rec}{par}
\olsection{Fungsi Rekursif Parsial}

Untuk memotivasi definisi fungsi rekursif, perhatikan bahwa bukti kita tentang
adanya fungsi yang dapat dihitung tetapi bukan rekursif primitif sebenarnya
menunjukkan hal yang jauh lebih kuat. Argumennya sederhana: kita hanya
menggunakan fakta bahwa fungsi-fungsi $f_0,f_1,\dots$ dapat dienumerasi
sedemikian sehingga $f_x(y)$, sebagai fungsi dari $x$ dan~$y$, dapat dihitung.
Jadi, argumen itu berlaku bagi \emph{setiap kelas fungsi yang dapat dienumerasi
dengan cara semacam itu}. Ini menempatkan kita dalam kesulitan: kita ingin
menguraikan fungsi-fungsi yang dapat dihitung secara eksplisit, tetapi setiap
uraian eksplisit tentang suatu kumpulan fungsi yang dapat dihitung tidak mungkin
mencakup semuanya!

Jalan keluarnya ialah memperbolehkan \emph{fungsi parsial}. Kita akan melihat
bahwa fungsi-fungsi parsial yang dapat dihitung \emph{memang} dapat dienumerasi.
\iftag{TMs}{ Sebenarnya, kita hampir sudah mengetahui bahwa demikianlah
  halnya, sebab mesin Turing dapat dienumerasi secara sistematis.}{} Kita akan
kembali ke argumen diagonal kita nanti dan menyelidiki mengapa argumen itu
tidak berhasil ketika fungsi parsial disertakan.

Sekarang pertanyaannya ialah: apa yang perlu kita tambahkan pada fungsi
rekursif primitif untuk memperoleh semua fungsi rekursif parsial? Kita perlu
melakukan dua hal:
\begin{enumerate}
\item Mengubah definisi fungsi rekursif primitif agar juga memperbolehkan
  fungsi parsial.
\item \emph{Menambahkan} sesuatu pada definisi tersebut agar beberapa fungsi
  parsial baru turut disertakan.
\end{enumerate}

Hal pertama mudah. Seperti sebelumnya, kita akan mulai dengan nol, penerus,
dan proyeksi, lalu menutup himpunan itu terhadap komposisi dan rekursi
primitif. Satu-satunya perbedaan ialah bahwa kita harus mengubah definisi
komposisi dan rekursi primitif untuk memperbolehkan kemungkinan bahwa beberapa
suku dalam definisi tidak terdefinisi. Jika $f$ dan $g$ merupakan fungsi
parsial, kita akan menulis $f(x) \fdefined$ untuk menyatakan bahwa $f$
terdefinisi pada~$x$, yakni $x$ berada dalam domain~$f$; dan $f(x)
\fundefined$ untuk menyatakan kebalikannya, yakni $f$ tidak terdefinisi
pada~$x$. Kita akan menggunakan $f(x)\simeq g(x)$ untuk menyatakan bahwa
$f(x)$ dan $g(x)$ keduanya tidak terdefinisi, atau keduanya terdefinisi dan
sama. Kita juga akan menggunakan notasi-notasi ini bagi suku yang lebih
rumit. Kita menerapkan konvensi bahwa jika $h$ dan $g_0$, \dots,~$g_k$
semuanya merupakan fungsi parsial, maka
\[
h(g_0(\vec x),\dots,g_k(\vec x))
\]
terdefinisi jika dan hanya jika setiap $g_i$ terdefinisi pada~$\vec x$ dan
$h$ terdefinisi ketika diterapkan pada nilai-nilai
$g_0(\vec x)$, \dots,~$g_k(\vec x)$. Dengan pengertian ini, definisi
komposisi dan rekursi primitif bagi fungsi parsial sama seperti sebelumnya,
kecuali bahwa kita harus mengganti ``$=$'' dengan ``$\simeq$''.

Yang akan kita tambahkan pada definisi fungsi rekursif primitif untuk
memperoleh fungsi parsial ialah \emph{operator pencarian tak berbatas}. Jika
$f(x,\vec z)$ merupakan sebarang fungsi parsial pada bilangan asli, definisikan
$\mu x\;f(x,\vec z)$ sebagai
\begin{quote}
  $x$ terkecil sedemikian sehingga $f(0,\vec z)$, $f(1,\vec z)$, \dots,
  $f(x,\vec z)$ semuanya terdefinisi dan $f(x,\vec z)=0$, jika $x$ semacam itu
  ada
\end{quote}
dengan pengertian bahwa $\mu x\;f(x,\vec z)$ tidak terdefinisi jika tidak ada.
Ini mendefinisikan $\mu x\;f(x,\vec z)$ secara unik.

\begin{explain}
Perhatikan bahwa definisi kita tidak merujuk pada\iftag{TMs}{ mesin Turing,
  atau}{} algoritme, ataupun model komputasi tertentu. Namun, seperti
komposisi dan rekursi primitif, pencarian tak berbatas didasari intuisi
operasional dan komputasional. Dalam komputabilitas fungsi parsial,
argumen-argumen yang padanya fungsi itu tidak terdefinisi bersesuaian dengan
masukan yang membuat
komputasi tidak berhenti. Prosedur untuk menghitung $\mu x\;f(x,\vec z)$
adalah sebagai berikut: hitung $f(0,\vec z)$, $f(1,\vec z)$,
$f(2,\vec z)$, dan seterusnya sampai nilai~$0$ dikembalikan. Akan tetapi,
jika salah satu komputasi perantara tidak berhenti, komputasi
$\mu x\;f(x,\vec z)$ juga tidak berhenti.
\end{explain}

Jika $R(x,\vec z)$ merupakan sebarang relasi, $\mu x\;R(x,\vec z)$
didefinisikan sebagai $\mu x\;(1\tsub\Char{R}(x,\vec z))$. Dengan kata lain,
$\mu x\;R(x,\vec z)$ mengembalikan nilai terkecil $x$ sedemikian sehingga
$R(x,\vec z)$ berlaku. Jadi, jika $f(x,\vec z)$ merupakan fungsi total,
$\mu x\;f(x,\vec z)$ sama dengan $\mu x\;(f(x,\vec z)=0)$. Namun,
perhatikan bahwa definisi asli kita lebih umum karena memperbolehkan
kemungkinan bahwa $f(x,\vec z)$ tidak terdefinisi untuk sebagian masukan (sedangkan,
sebaliknya, fungsi karakteristik suatu relasi selalu total).

\begin{defn}
Himpunan \emph{fungsi rekursif parsial} ialah himpunan terkecil fungsi parsial
pada bilangan asli (dengan berbagai aritas) yang memuat nol, penerus, dan
proyeksi, serta tertutup terhadap komposisi, rekursi primitif, dan pencarian
tak berbatas.
\end{defn}

Tentu saja, beberapa fungsi rekursif parsial kebetulan bersifat total, yakni
terdefinisi bagi setiap argumen.

\begin{defn}
\ollabel{defn:recursive-fn}
Himpunan \emph{fungsi rekursif} ialah himpunan fungsi rekursif parsial yang
bersifat total.
\end{defn}

Suatu fungsi rekursif kadang-kadang disebut ``rekursif total'' untuk
menegaskan bahwa fungsi itu terdefinisi pada setiap masukan.

\end{document}
